\chapter{Fundamental Theorem of Calculus}

\section{Convexity}

\begin{definition}[Raw-real order]\label{def:raw-real-order}
\lean{ComputableAnalysis.RealRaw.Le}
\uses{def:raw-real}
\leanok
For raw reals \(X,Y\), define
\[
  X\le Y
\]
to mean that every rational lower approximation to \(X\) is at most every
rational upper approximation to \(Y\):
\[
  \forall n,m,\qquad X_n^- \le Y_m^+.
\]
This is the order relation used in mathematical statements.  For valid raw
representatives,
\[
  X\sim Y
  \quad\Longleftrightarrow\quad
  X\le Y\ \text{and}\ Y\le X.
\]
\end{definition}

\begin{definition}[Endpoint difference]\label{def:endpoint-diff}
\lean{ComputableAnalysis.endpointDifferenceInterval}
\uses{def:real-fun-raw,def:raw-real-arithmetic}
For a raw real-valued function \(F\), define
\[
  E(F;u,v)=F(v)-F(u).
\]
At stage \(p\), if \(F(u)_p=[A^-,A^+]\) and
\(F(v)_p=[B^-,B^+]\), then
\[
  E(F;u,v)_p=[B^- - A^+,B^+ - A^-].
\]
\end{definition}

\begin{definition}[Rational secant]\label{def:rational-secant}
\lean{ComputableAnalysis.secantRaw}
\uses{def:endpoint-diff,def:raw-real-order}
\leanok
For rational \(u<v\), the secant of \(F\) on \([u,v]\) is the raw real
\[
  \Sec_F(u,v)=\frac{F(v)-F(u)}{v-u}.
\]
Its finite stages are the rational interval slopes
\[
  \Sec_F(u,v)_p=E(F;u,v)_p/(v-u).
\]
\end{definition}

\begin{definition}[Convexity]\label{def:exact-convexity}
\lean{ComputableAnalysis.ExactConvexOn}
\uses{def:rational-secant,def:raw-real-order}
\leanok
A raw real-valued function \(F\) is convex on the rational interval
\([a,b]\) when it is defined and valid at every rational point of
\([a,b]\), and its rational secants are monotone:
\[
  \Sec_F(w,x)\le \Sec_F(y,z)
\]
whenever
\[
  a\le w<x\le y<z\le b.
\]
Concavity is the same statement with the secant order reversed.
\end{definition}

\section{Derivative of a Convex Function}

\begin{definition}[Right derivative of a convex function]
\label{def:convex-right-derivative}
\lean{ComputableAnalysis.RightDerivativeAt}
\uses{def:exact-convexity}
Let \(F\) be convex on \([a,b]\), and let \(q<b\) be rational.  The right
secants
\[
  \Sec_F(q,q+h),\qquad h>0,\quad q+h\le b,
\]
are monotone as \(h\downarrow 0\) and bounded below by any left secant ending
at \(q\).  The right derivative \(D_+F(q)\) is represented by the infimum of
these right secants:
\[
  D_+F(q)=\inf_{0<h,\ q+h\le b}\Sec_F(q,q+h).
\]
Equivalently, it is the limit of any rational sequence \(h_n\downarrow0\)
with \(q+h_n\le b\).
\end{definition}

\begin{definition}[Left derivative of a convex function]
\label{def:convex-left-derivative}
\lean{ComputableAnalysis.LeftDerivativeAt}
\uses{def:exact-convexity}
For \(a<q\), the left derivative is
\[
  D_-F(q)=\sup_{0<h,\ a\le q-h}\Sec_F(q-h,q).
\]
For convex \(F\), \(D_-F(q)\le D_+F(q)\).  The ordinary two-sided derivative
exists exactly at the rational points where these two one-sided derivatives
are equivalent.
\end{definition}

\begin{theorem}[Right derivative is monotone]
\label{thm:convex-right-derivative-increasing}
\lean{ComputableAnalysis.rightDerivativeAt_mono}
\uses{def:convex-right-derivative}
\leanok
If \(F\) is convex on \([a,b]\), then
\[
  q_1<q_2
  \quad\Longrightarrow\quad
  D_+F(q_1)\le D_+F(q_2)
\]
for rational \(q_1,q_2<b\).
\end{theorem}

\begin{proof}
Choose rational \(h_1,h_2>0\) so small that
\[
  q_1+h_1\le q_2\le q_2+h_2.
\]
Convexity gives
\[
  \Sec_F(q_1,q_1+h_1)\le \Sec_F(q_2,q_2+h_2).
\]
Taking the infimum over right secants on both sides gives the claim.
\end{proof}

\begin{definition}[Derivative formula identification]\label{def:formula-identification}
\uses{def:convex-right-derivative}
Let \(g\) be a proposed raw function formula.  On a rational interval,
\(g\) is identified with \(D_+F\) when it is equivalent to the right-secant
infimum defining \(D_+F\).  The formula is not a guess: it is proved from
finite secant inequalities or algebraic identities for the particular
function.
\end{definition}

\begin{example}[Trigonometric functions]\label{ex:trig-derivatives}
\uses{def:formula-identification,thm:geometric-trig-identities,def:geometric-trig-integral-identities}
For the geometric sine and cosine from the circle chapter, convexity or
concavity is checked on intervals where the sign of the appropriate second
derivative is certified by the circle algorithm.  The derivative formulas are
\[
  (\sin)'=\cos,\qquad (\cos)'=-\sin.
\]
The integral identities in Chapter~\ref{ch:integrals} are the corresponding
definite-integral statements once these derivative certificates are supplied.
\end{example}

\begin{example}[Exponential]\label{ex:exp-derivative}
\uses{def:formula-identification,def:exp-reps}
For the exponential functions constructed in the exponential and logarithm
chapter, convexity follows from positivity and the derivative formula is
\[
  (\exp)'=\exp.
\]
The power-series construction proves this by termwise differentiation; the
Euler-limit construction must then be compared with the same derivative
formula through the representation-agreement theorem.
\end{example}

\begin{example}[Logarithm]\label{ex:log-ftc}
\uses{def:formula-identification,def:log-integral}
For logarithm on \(0<a\le x\le b\), the derivative formula is
\[
  (\log)'(x)=\frac1x.
\]
The positivity certificate supplies the domain-apartness data needed for
interval division, and the secant inequalities identify \(1/x\) with the
one-sided derivative.  Since \(1/x\) is decreasing on positive intervals,
\(\log\) is concave there.
\end{example}

\begin{example}[Geometric arctangent]\label{ex:arctan-ftc}
\uses{def:formula-identification,def:inverse-elementary-integral-identities}
Let \(A(x)\) be the geometric arctangent branch.  On \(0\le u<v\), the finite
sector-area inequality
\[
  \frac{1}{1+v^2}
  \le
  \Sec_A(u,v)
  \le
  \frac{1}{1+u^2}
\]
identifies the derivative with \(1/(1+x^2)\).  This is the bridge from the
geometric construction to the arctangent integral and Taylor-series
representations.
\end{example}

\begin{remark}[Corners]\label{rem:convex-corners}
\uses{def:convex-left-derivative,def:convex-right-derivative}
Convexity does not imply existence of the two-sided derivative.  The function
\(|x|\) is convex and has \(D_-F(0)<D_+F(0)\).  The universal FTC for convex
functions must therefore use a one-sided derivative, or else explicitly
restrict to the points where the left and right derivatives agree.
\end{remark}

\section{Fundamental Theorem of Calculus}

The integral used in this chapter is the constructive definite integral
already introduced in Chapter~\ref{ch:integrals}.  The present chapter proves
when such an integral is equal to an endpoint difference.

\begin{theorem}[Monotone functions are integrable]
\label{thm:monotone-integrable}
\uses{def:definite-integral-identity,def:raw-real-order}
Every monotone rational-input raw real-valued function on a rational interval
has a constructive definite integral.  Monotonicity supplies the lower and
upper cell values from endpoint evaluations, and Darboux sums converge under
rational refinement.
\end{theorem}

\begin{theorem}[Fundamental theorem of calculus]\label{thm:convex-ftc}
\uses{def:exact-convexity,def:convex-right-derivative,thm:monotone-integrable,def:definite-integral-identity}
If \(F\) is convex on a rational interval \([a,b]\), then the right
derivative \(D_+F\) is integrable and
\[
  \int_a^b D_+F(x)\,\dd x \sim F(b)-F(a).
\]
Equivalently, the same conclusion holds with the left derivative \(D_-F\).
\end{theorem}

\begin{proof}
For a partition \(P:a=x_0<\cdots<x_N=b\), convexity gives
\[
  D_+F(x_i)
  \le
  \Sec_F(x_i,x_{i+1})
  \le
  D_-F(x_{i+1})
\]
on each cell.  Multiplying by \(x_{i+1}-x_i\) and summing yields Darboux
lower and upper sums enclosing
\[
  \sum_i \bigl(F(x_{i+1})-F(x_i)\bigr)=F(b)-F(a).
\]
The gap between the sums is controlled by monotonicity of the one-sided
derivatives, and tends to zero under rational refinement.  Hence the integral
and the endpoint difference represent the same raw real.
\end{proof}

\begin{theorem}[Formula identification closure]
\label{thm:formula-identification-ftc}
\uses{def:formula-identification,thm:convex-ftc}
If \(g\) is identified with \(D_+F\) on \([a,b]\), then
\[
  \int_a^b g(x)\,\dd x \sim F(b)-F(a).
\]
\end{theorem}

\begin{remark}[Piecewise convexity]\label{rem:no-piecewise-ftc}
\uses{thm:convex-ftc}
There is no separate named theorem for piecewise convex or concave functions.
In an example proof, split the interval at rational breakpoints where the
convexity proof changes, apply \cref{thm:convex-ftc} on each piece, and
combine the endpoint equalities by raw-real arithmetic and transitivity.
\end{remark}

\section{Taylor's Theorem}

\begin{definition}[Taylor by iterated FTC]\label{def:taylor-iterated-ftc}
\lean{ComputableAnalysis.Taylor.IteratedFTCChain}
\uses{thm:convex-ftc,thm:formula-identification-ftc}
A Taylor expansion is generated by repeatedly applying the FTC on rational
subintervals where the needed convexity or concavity certificates are
available:
\[
  F(x)
  =
  F(a)+\int_a^x F'(t)\,\dd t.
\]
Iterating this statement gives the polynomial part plus an explicit
iterated-integral remainder.
\end{definition}

\begin{theorem}[Taylor's theorem with integral remainder]\label{thm:taylor-integral-remainder}
\uses{def:taylor-iterated-ftc}
If \(F\) has the required derivative-identification certificates through
order \(n+1\) on a rational interval, then
\[
  F(x)
  =
  \sum_{k=0}^{n} \frac{F^{(k)}(a)}{k!}(x-a)^k
  +
  R_n(x),
\]
where the remainder is represented by an iterated constructive integral of
\(F^{(n+1)}\).  Bounding that final integral gives the usual effective
Taylor estimate.
\end{theorem}

\begin{theorem}[Arctangent kernel remainder]\label{thm:atan-kernel-rem}
\lean{ComputableAnalysis.Taylor.ArctanKernel.finiteRemainderRoute}
\uses{def:taylor-iterated-ftc}
\leanok
For rational \(x\), the algebraic identity
\[
  \frac{1}{1+x^2}
  =
  \sum_{k=0}^{n} (-1)^k x^{2k}
  +
  \frac{(-1)^{n+1}x^{2n+2}}{1+x^2}
\]
holds with the effective bound
\[
  \left|
    \frac{(-1)^{n+1}x^{2n+2}}{1+x^2}
  \right|
  \le |x|^{2n+2}.
\]
This is the finite algebra behind the arctangent Taylor remainder.
\end{theorem}
